% !TeX root = ../教材题库_全部条目_转写工作稿.tex
% Source: XJ-B1-C04-HW-P126-Y0578
% 第4章 / 4.3 对数函数 / 教材 p.126 / PDF 第133页
% OCR未校对：公式与文字均可能有误；未完成数学TeX转写；禁止直接发给学生
% 原题及相邻原解截图（核对用）：../题目原图/XJ-B1-C04-HW-P126-Y0578-01.webp
% 本文件只在转写完成后作为题干使用；OCR 条目尚不是可靠题干。
\begin{bankitem}{习题4.3 第17题}{第4章 / 4.3 对数函数 / 教材 p.126 / PDF 第133页}{XJ-B1-C04-HW-P126-Y0578}{OCR未校对：公式与文字均可能有误；未完成数学TeX转写；禁止直接发给学生}
\begin{ocrdraft}
[教材p.126 / PDF第133页]
\par
17．记f(x)=ax，g(x)=logax(a>()且乒1），写出f(g(x))和g(f(x))的具
\par
体表达式并化简．尝试说明这里的表达式和化简结果与对数的基本恒等式
\par
和b=logaab的关系．
\end{ocrdraft}
\end{bankitem}
